\chapter{A Economia Sacramental}

A \textbf{economia sacramental} é a maneira pela qual Deus torna presente, no tempo e na história, a salvação realizada por Cristo. O Catecismo ensina que o plano de Deus --- a ``economia da salvação'' --- se realiza de modo culminante no \textbf{mistério pascal}, e que a \textbf{liturgia} (sobretudo os sacramentos e, de modo particular, a Eucaristia) faz com que os fiéis \emph{vivam desse mistério} e o expressem na vida e deem testemunho no mundo.

\section*{1066--1075: A economia da salvação e o mistério pascal celebrado na liturgia}
No \textbf{Símbolo da fé}, a Igreja confessa o plano do ``bom prazer'' de Deus para toda a criação: o Pai realiza o ``mistério de sua vontade'' ao dar o Filho amado e o Espírito Santo para a salvação do mundo e para a glória do nome de Deus.

Esse desígnio é apresentado como a \textbf{misteriosa iniciativa de Deus no tempo}: a tradição chama esse conjunto ordenado de ações divinas de \textbf{``economia da Palavra encarnada''} ou de \textbf{``economia da salvação''}.

Em continuidade com o Antigo Testamento, os ``maravilhosos feitos de Deus'' entre o povo eram apenas um \textbf{anúncio} do que Cristo realizaria para redimir a humanidade e dar glória perfeita a Deus; Cristo cumpriu essa obra \textbf{principalmente} pelo \textbf{mistério pascal} da Paixão, Ressurreição e Ascensão: ``morrendo, destruiu a nossa morte, rising he restored our life.''

Por isso, a Igreja celebra na liturgia \textbf{sobretudo} o mistério pascal, pois é nesse mesmo mistério de Cristo que a Igreja \textbf{proclama e celebra} aquilo que permite aos fiéis \emph{viverem} dele e \emph{testemunharem} no mundo.

O Catecismo afirma explicitamente que ``na liturgia, especialmente no sacrifício divino da Eucaristia, se realiza a obra da nossa redenção''; e é através da liturgia que os fiéis se tornam capazes de ``expressar em suas vidas e manifestar aos outros'' o mistério de Cristo e a verdadeira natureza da Igreja.

\emph{(Observação: os parágrafos 1069--1075 não aparecem nas fontes fornecidas nesta conversa; portanto, não é possível completar integralmente essa faixa com fidelidade textual.)}

\section**{1077--1112: A liturgia como bênção divina e resposta humana}
Nesta seção, as fontes fornecidas trazem o Catecismo trabalhando o conceito de \textbf{bênção} no horizonte da ação de Deus.

O Catecismo recorda a atitude de Deus que ``nos abençoou em Cristo com toda bênção espiritual... escolheu-nos nele... para sermos santos e irrepreensíveis''.

A bênção, segundo o Catecismo, é uma ação divina e vivificante, cuja fonte é o Pai; é ao mesmo tempo \textbf{palavra} e \textbf{dom}. Quando aplicada ao ser humano, a palavra ``bênção'' significa \textbf{adoração} e \textbf{entrega} ao Criador em \textbf{ação de graças}.

Desde o começo até o fim do tempo, ``toda a obra de Deus é bênção''; dos cânticos do primeiro ato da criação até os hinos da Jerusalém celeste, os autores inspirados apresentam o plano da salvação como uma imensa bênção divina.

\emph{(Observação: os parágrafos 1080--1112 não aparecem nas fontes fornecidas nesta conversa; portanto, não é possível completar integralmente essa faixa com fidelidade textual.)}

\section*{1113--1134: O que é comum aos sacramentos na economia sacramental}
A catequese dos sacramentos ensina que a vida litúrgica da Igreja gira em torno do \textbf{sacrifício eucarístico} e dos \textbf{sacramentos}, e que há \textbf{sete sacramentos} na Igreja: Batismo, Confirmação (ou Crisma), Eucaristia, Penitência, Unção dos Enfermos, Ordem e Matrimônio.

O Catecismo afirma que os ``sacramentos da nova lei'' foram \textbf{instituídos por Jesus Cristo, nosso Senhor}. Essa fonte comum sublinha a unidade da economia sacramental: os sacramentos não são invenções posteriores, mas dispensações do próprio Cristo à sua Igreja.

O Catecismo explica ainda que as \textbf{palavras e ações de Jesus} durante a vida oculta e o ministério público já eram \textbf{salvíficas}, pois \textbf{antecipavam} a força do mistério pascal; elas anunciavam e preparavam aquilo que Cristo iria dar à Igreja quando tudo estivesse cumprido.

Por isso, ``os mistérios da vida de Cristo'' são o fundamento do que, doravante, Ele ``dispensará'' nos sacramentos, por meio dos ministros da Igreja: ``o que era visível em nosso Salvador passou para os seus mistérios''.

Além disso, os sacramentos são descritos como ``\textbf{poderes que provêm}'' do Corpo de Cristo, que é sempre vivo e vivificante; são ações do Espírito Santo que atua no seu Corpo, a Igreja; e são as ``\textbf{obras-primas de Deus}'' na nova e eterna aliança.

\emph{(Observação: os parágrafos 1116--1134 não aparecem nas fontes fornecidas nesta conversa; portanto, não é possível completar integralmente essa faixa com fidelidade textual.)}

\section*{1135--1144: Quem celebra e como a liturgia é celebrada --- sinais e símbolos}
O Catecismo inicia esta parte situando a catequese da liturgia: ela pressupõe primeiro compreender a \textbf{economia sacramental} (Capítulo Um). Nessa luz, revela-se a ``inovação'' própria de sua celebração: o conteúdo passa a ser tratado como celebração dos sacramentos da Igreja, comuns ao conjunto dos sete, e em seguida o que é próprio de cada um.

O Catecismo também organiza as perguntas fundamentais da catequese litúrgica: \textbf{quem celebra}, \textbf{como}, \textbf{quando} e \textbf{onde} é celebrada a liturgia.

Quanto a \textbf{quem celebra}, a liturgia é apresentada como ``\textbf{ação} do \emph{Cristo total} (Christus totus)''; aqueles que a celebram sem sinais ainda assim já participam da liturgia celeste, onde celebrar é comunhão e festa.

O Catecismo recorre à visão do Apocalipse para apresentar os celebrantes da liturgia celeste: um trono no céu (Deus), o Cordeiro como imolado, Cristo crucificado e ressuscitado, sacerdote do verdadeiro santuário, e por fim o ``rio de água da vida'' que simboliza o Espírito Santo.

O Catecismo ensina que é ``na liturgia eterna'' que o Espírito e a Igreja nos fazem participar sempre que celebramos o mistério da salvação nos sacramentos.

Quanto a \textbf{como} a liturgia é celebrada, o Catecismo sublinha que ``uma celebração sacramental é tecida de sinais e símbolos''. Esses sinais, conforme a ``pedagogia divina da salvação'', têm seu sentido enraizado na criação e na cultura humana, especificado pelos eventos da Antiga Aliança e revelado plenamente na pessoa e obra de Cristo.

\emph{(Observação: os parágrafos 1138--1144 não aparecem nas fontes fornecidas nesta conversa; portanto, não é possível completar integralmente essa faixa com fidelidade textual.)}

\section*{1145--1162: Sinais do mundo humano e linguagem da criação}
\subsection*{1145--1147: A pedagogia dos sinais}
O Catecismo afirma: ``A celebração sacramental é tecida de sinais e símbolos''; seus significados são enraizados na obra da criação e na cultura humana, especificados pelos eventos da Antiga Aliança e plenamente revelados na pessoa e obra de Cristo.

Os sinais têm papel decisivo na vida humana: como o ser humano é ao mesmo tempo corpo e espírito, expressa e percebe realidades espirituais por meios físicos; como ser social, necessita de sinais e símbolos para comunicar-se (linguagem, gestos e ações); e isso vale também para sua relação com Deus.

Nesse mesmo horizonte, o Catecismo ensina que \textbf{Deus fala ao ser humano por meio da criação visível}: o cosmos material é apresentado à inteligência humana de modo que nela se possam ler vestígios do Criador.

E exemplos concretos ajudam a entender essa linguagem: luz e trevas, vento e fogo, água e terra, a árvore e seu fruto falam de Deus e simbolizam ao mesmo tempo sua grandeza e sua proximidade.

\emph{(Observação: os parágrafos 1148--1162 não aparecem nas fontes fornecidas nesta conversa; portanto, não é possível completar integralmente essa faixa com fidelidade textual.)}

\section*{1163--1178: O ano litúrgico --- o ``Hoje!'' de Deus e o desdobramento do mistério de Cristo}
O Catecismo apresenta a prática do \textbf{ano litúrgico} como comemoração da obra salvífica do ``divino Esposo'' da Igreja. Ao longo do ano, a Igreja celebra o mistério de Cristo: semanalmente, no Dia do Senhor, lembra a ressurreição; anualmente, no tempo de Páscoa, celebra com especial solenidade a Paixão; e, no curso do ano, ``desdobra todo o mistério de Cristo''.

Ao fazê-lo, a Igreja ``reconhece o rememorado'': recordando os mistérios da redenção, abre aos fiéis as riquezas das potências e méritos do Senhor, tornando-os de algum modo presentes em todas as épocas; assim, os fiéis se apoderam deles e são preenchidos com graça salvadora.

O Catecismo explica também que, desde a Lei de Moisés, o povo de Deus observava festas fixas para comemorar ações extraordinárias do Salvador, agradecer e perpetuar o memorial, formando novas gerações. Na idade da Igreja, entre a Páscoa de Cristo já realizada e a consumação no Reino de Deus, a liturgia celebrada em dias fixos traz a marca da novidade do mistério de Cristo.

Quando a Igreja celebra o mistério de Cristo, surge uma palavra que marca sua oração: \textbf{``Hoje!''}. Esse ``hoje'' do Deus vivo que o homem é chamado a entrar é \textbf{``a hora'' do Passo (Passover) de Jesus}, que atravessa e sustenta toda a história.

Assim, o Catecismo descreve a experiência do crente como abertura ao dia de luz duradoura, a \textbf{Páscoa mística}, que não é apagada para aqueles que creem.

\emph{(Observação: as fontes fornecidas terminam no parágrafo 1165; os parágrafos 1166--1178 não aparecem nesta conversa; portanto, não é possível completar integralmente essa faixa.)}

\section*{1179--1199 e 1200--1209}
As fontes fornecidas nesta conversa não incluem os textos dos parágrafos \textbf{1179--1199} nem \textbf{1200--1209}. Assim, não é possível reformular essas faixas de modo fiel a partir do material recebido.

\section*{Síntese final}
Com base nas passagens disponíveis, a economia sacramental aparece como a \textbf{atualização no tempo} da salvação: Deus realiza o seu plano (economia da salvação) culminando no mistério pascal; a Igreja proclama e celebra esse mistério na liturgia para que os fiéis vivam dele e deem testemunho; e a liturgia educa e comunica por meio da ação sacramental, isto é, por \textbf{sinais} e por uma \textbf{linguagem} que brota da criação e da cultura, oferecendo ao homem o ``Hoje!'' de Deus na história do mistério de Cristo.